\documentclass[twoside,spanish,a4paper,12pt]{tfg}

\titulacion{Grado en Ingeniería Informática}
\title{Qsimov: Re-entrenamiento Incremental sin Olvido Catastrófico en Redes Neuronales con Activación ReLU}
\authorlabel{Autor}
\author{Matías Moncho Ronda}
\tutorlabel{Tutor}
\tutor{Dr. Juan Manuel Orduña\\\hspace*{5cm}Huertas}
\convocatoria{Julio 2026}

\usepackage{booktabs}
\usepackage{multirow}
\usepackage{url}
\usepackage[protrusion=true,expansion=false]{microtype}
\setlength{\emergencystretch}{60pt}
\DeclareUrlCommand\bpath{\urlstyle{tt}\def\UrlFont{\bfseries\ttfamily}}

\AtBeginDocument{\shorthandoff{"}\catcode`\'=12}

\begin{document}

\portada
\cleardoublepage
\contraportada
\cleardoublepage
\ia
\cleardoublepage

% ---- Resumen en español ----
\begin{resumen}
El aprendizaje continuo en redes neuronales artificiales choca con un problema estructural conocido
como olvido catastrófico: cuando una red es re-entrenada con nuevas muestras, el ajuste de sus pesos
destruye el conocimiento previamente adquirido. Este trabajo evalúa \textbf{Qsimov}, una librería
Python existente que aborda este problema mediante la explotación de la estructura lineal por tramos
que inducen las activaciones ReLU en redes feed-forward.

La idea central consiste en representar el comportamiento de las últimas capas de una red
pre-entrenada en términos de sus \textbf{caminos activos}: las combinaciones de conexiones entre
neuronas que permanecen activas para cada muestra dada. Esta representación transforma el
re-entrenamiento en la resolución de un sistema de ecuaciones lineales cuyas incógnitas son los pesos
de los caminos, y no los pesos individuales de la red original. Gracias a esta reformulación, es
posible acumular las ecuaciones de distintas rondas de entrenamiento sin necesidad de acceder a los
datos anteriores, eliminando así el olvido.

La librería ofrece dos mecanismos de re-entrenamiento: resolución directa mediante el sistema lineal
(variante más rápida) y descenso de gradiente sobre una capa completamente conectada restringida a
los caminos activos (variante más flexible). Ambos se implementan sobre TensorFlow/Keras y PyTorch
con extensiones en Cython para las operaciones más costosas.

Los experimentos, realizados sobre MNIST, CIFAR-10 e ImageNet con 100 clases, demuestran que el
sistema lineal con acumulación de ecuaciones elimina el olvido catastrófico en todos los escenarios
evaluados: mantiene 94.3\,\% de precisión en clases antiguas en MNIST (frente al 0.0\,\% del
fine-tuning estándar), 73.4\,\% en CIFAR-10, y entre el 79\,\% y el 84\,\% en rondas previas sobre
ImageNet, todo ello entre seis y diez veces más rápido que el re-entrenamiento acumulado estándar.

\end{resumen}
\cleardoublepage

% ---- Abstract en inglés ----
\begin{abstract}
Continual learning in artificial neural networks is hindered by a structural problem known as
catastrophic forgetting: when a network is retrained on new samples, the weight updates destroy
previously acquired knowledge. This work evaluates \textbf{Qsimov}, an existing Python library that
addresses this problem by exploiting the piecewise-linear structure induced by ReLU activations in
feed-forward networks.

The core idea is to represent the behaviour of the final layers of a pre-trained network in terms of
their \textbf{active paths}: the combinations of neuron connections that remain active for each given
sample. This representation transforms retraining into solving a linear system of equations whose
unknowns are the path weights, not the individual weights of the original network. Thanks to this
reformulation, equations from different training rounds can be accumulated without requiring access
to previous data, thereby eliminating forgetting.

The library offers two retraining mechanisms: direct solving via the linear system (faster variant)
and gradient descent over a fully connected layer restricted to the active paths (more flexible
variant). Both are implemented for TensorFlow/Keras and PyTorch, with Cython extensions for the most
computationally intensive operations.

Experiments conducted on MNIST, CIFAR-10 and ImageNet with 100 classes demonstrate that the linear
system with equation accumulation eliminates catastrophic forgetting in all evaluated scenarios: it
maintains 94.3\,\% accuracy on old classes in MNIST (versus 0.0\,\% for standard fine-tuning),
73.4\,\% in CIFAR-10, and between 79\,\% and 84\,\% in previous rounds on ImageNet, all between
six and ten times faster than standard cumulative retraining.

\end{abstract}
\cleardoublepage

% ---- Resum en valencià ----
\begin{resum}
L'aprenentatge continu en xarxes neuronals artificials es troba amb un problema estructural conegut
com a oblit catastròfic: quan una xarxa és re-entrenada amb noves mostres, l'ajust dels seus pesos
destrueix el coneixement prèviament adquirit. Aquest treball avalua \textbf{Qsimov}, una llibreria
Python existent que aborda aquest problema mitjançant l'explotació de l'estructura lineal a trams que
indueixen les activacions ReLU en xarxes feed-forward.

La idea central consisteix a representar el comportament de les últimes capes d'una xarxa
pre-entrenada en termes dels seus \textbf{camins actius}: les combinacions de connexions entre
neurones que romanen actives per a cada mostra donada. Aquesta representació transforma el
re-entrenament en la resolució d'un sistema d'equacions lineals les incògnites de les quals són els
pesos dels camins, i no els pesos individuals de la xarxa original. Gràcies a aquesta reformulació,
és possible acumular les equacions de distintes rondes d'entrenament sense necessitat d'accedir a les
dades anteriors, eliminant així l'oblit.

La llibreria ofereix dos mecanismes de re-entrenament: resolució directa mitjançant el sistema lineal
(variant més ràpida) i descens del gradient sobre una capa completament connectada restringida als
camins actius (variant més flexible). Tots dos s'implementen sobre TensorFlow/Keras i PyTorch amb
extensions en Cython per a les operacions més costoses.

Els experiments, realitzats sobre MNIST, CIFAR-10 i ImageNet amb 100 classes, demostren que el
sistema lineal amb acumulació d'equacions elimina l'oblit catastròfic en tots els escenaris avaluats:
manté un 94.3\,\% de precisió en classes antigues en MNIST (enfront del 0.0\,\% del fine-tuning
estàndard), un 73.4\,\% en CIFAR-10, i entre el 79\,\% i el 84\,\% en rondes prèvies sobre
ImageNet, tot plegat entre sis i deu vegades més ràpid que el re-entrenament acumulat estàndard.

\end{resum}
\cleardoublepage

\tableofcontents

\pagestyle{tfg}
\justify

\chapter{Introducción y Motivación}
% Capítulo 1: Introducción y Motivación
% Cubre las secciones 1 y 2 de la memoria

\section{Introducción}

Las redes neuronales profundas dominan actualmente el estado del arte en tareas de visión por
computador, procesamiento del lenguaje natural y reconocimiento de patrones. Sin embargo, el
paradigma dominante de entrenamiento (procesar todos los datos de una vez con acceso
irrestricto) resulta inapropiado en muchos escenarios reales. Los datos llegan de forma
secuencial, los recursos computacionales son limitados, y las restricciones de privacidad o
almacenamiento impiden conservar todos los datos históricos. En este contexto surge el campo del
\textbf{aprendizaje continuo} (\textit{continual learning}), cuyo objetivo es que un modelo aprenda
de nuevos datos sin olvidar los conocimientos previos.

El obstáculo principal de este campo es el \textbf{olvido catastrófico}
(\textit{catastrophic forgetting}), descrito por primera vez de forma sistemática por McCloskey y
Cohen en 1989~\cite{ref1}. Cuando una red neuronal es re-entrenada sobre nuevas muestras mediante
descenso de gradiente, la actualización de sus pesos interfiere destructivamente con las
representaciones aprendidas anteriormente, de forma mucho más severa que el olvido observado en
sistemas biológicos.

La mayoría de las soluciones propuestas en la literatura (regularización, expansión de
arquitectura, o replay de datos generados) introducen compromisos entre plasticidad (capacidad de
aprender cosas nuevas) y estabilidad (capacidad de retener las anteriores), y frecuentemente
requieren almacenar ejemplos pasados o aumentar la complejidad del modelo con cada nueva tarea.

Este trabajo evalúa Qsimov, una librería que implementa un enfoque alternativo basado en una
propiedad matemática de las redes con activación ReLU: para cualquier entrada, la función que
computa la red es \textbf{lineal} en la región del espacio que contiene a esa entrada. Cada entrada
traza un <<camino>> único a través de la red, determinado por qué neuronas están activas ($> 0$) y
cuáles no. Qsimov separa la estructura de esos caminos, aprendida durante el pre-entrenamiento,
de los pesos asociados a cada camino, que se resuelven analíticamente mediante un sistema de
ecuaciones lineales. La acumulación incremental de ecuaciones a lo largo del tiempo garantiza que el
sistema <<recuerde>> todas las rondas de entrenamiento sin necesidad de acceder a los datos
originales.

El presente documento describe el funcionamiento de Qsimov y los experimentos realizados para
evaluar su comportamiento en escenarios de aprendizaje continual, siguiendo la estructura
recomendada para trabajos de fin de grado de ingeniería en informática~\cite{ref2}.

\section{Motivación y Objetivos}

\subsection{Motivación}

La motivación del proyecto nace de una observación fundamental sobre la naturaleza del entrenamiento
de redes neuronales con activación ReLU. Una red feed-forward con ReLU define, implícitamente, una
partición del espacio de entrada en regiones lineales: dentro de cada región, la función que computa
la red es una transformación lineal afín~\cite{ref3}. Esa transformación queda completamente
determinada por el subconjunto de neuronas que están activas para las entradas de esa región, es
decir, por el <<camino activo>> de la muestra.

Esta propiedad tiene una consecuencia directa para el re-entrenamiento: si se fijan los parámetros
que determinan \textit{qué} caminos existen (los pesos del modelo base, aprendidos durante el
pre-entrenamiento), el problema de encontrar los pesos óptimos para esos caminos se reduce a
resolver un sistema de ecuaciones lineales. Y un sistema lineal puede ampliarse con nuevas
ecuaciones en cualquier momento, incorporando nuevas muestras o nuevas clases, sin que las
soluciones previas se destruyan, siempre que las ecuaciones antiguas permanezcan en el sistema.

Frente a las soluciones existentes basadas en penalizaciones (EWC~\cite{ref4}, SI~\cite{ref5}) o
en expansión de arquitectura (PNN~\cite{ref6}), Qsimov ofrece garantías exactas de no-olvido para
las muestras incluidas en el sistema, sin requerir hiperparámetros adicionales y sin aumentar el
tamaño de la red.

\subsection{Objetivos}

Los objetivos del proyecto son los siguientes:

El objetivo principal del proyecto es evaluar experimentalmente la librería Qsimov en escenarios de
aprendizaje continual, caracterizando su comportamiento frente al olvido catastrófico y
comparándolo con métodos de re-entrenamiento estándar.

Los objetivos secundarios son:

\begin{enumerate}
  \item Comprender el funcionamiento de los componentes principales de Qsimov
        (\texttt{PathSelector}, \texttt{QsimovLinearSystem}, \texttt{QsimovGradient}) y su API de
        uso.
  \item Diseñar y ejecutar experimentos de evaluación sobre tres \textit{benchmarks} de distinta
        escala y complejidad: MNIST, CIFAR-10 e ImageNet (100 clases).
  \item Comparar el rendimiento de Qsimov (variante de sistema lineal y variante de gradiente) con
        métodos de referencia: fine-tuning estándar y re-entrenamiento acumulado.
  \item Analizar el impacto del punto de corte (\texttt{initial\_layer}) sobre la precisión
        obtenida y el número de caminos generados.
  \item Evaluar la robustez al learning rate del módulo de gradiente frente al re-entrenamiento
        estándar.
  \item Identificar los escenarios en los que Qsimov ofrece ventajas claras y aquellos en los que
        presenta limitaciones.
\end{enumerate}


\chapter{Estado del Arte}
% Capítulo 2: Estado del Arte

\section{El problema del olvido catastrófico}

El olvido catastrófico fue caracterizado por primera vez de forma rigurosa por McCloskey y
Cohen~\cite{ref1} y Ratcliff~\cite{ref7}, que demostraron experimentalmente cómo incluso un
entrenamiento mínimo sobre nueva información borraba casi por completo las representaciones
previamente aprendidas en modelos conexionistas. A diferencia del olvido humano, que es gradual y
selectivo, el olvido en redes artificiales es abrupto y masivo: basta con actualizar los pesos
mediante descenso de gradiente sobre un nuevo conjunto de datos para que el rendimiento en los datos
anteriores caiga drásticamente.

La raíz del problema reside en la naturaleza distribuida de las representaciones en redes
neuronales: los mismos pesos participan en la codificación de múltiples patrones, de modo que su
modificación para adaptarse a los nuevos datos necesariamente interfiere con los antiguos. Este
conflicto entre capacidad de aprender cosas nuevas y capacidad de retener las ya aprendidas se
denomina el \textbf{dilema estabilidad-plasticidad}~\cite{ref8}.

\section{Taxonomía de soluciones}

Las estrategias para mitigar el olvido catastrófico se agrupan en tres grandes familias~\cite{ref9}:

\subsection{Métodos basados en regularización}

Estos métodos añaden un término de penalización a la función de pérdida para evitar que los pesos
más importantes para las tareas anteriores cambien demasiado. El más influyente es \textbf{EWC}
(\textit{Elastic Weight Consolidation}, Kirkpatrick et al., 2017)~\cite{ref4}, que usa la diagonal
de la matriz de información de Fisher como medida de la importancia de cada peso. Formalmente, la
función de pérdida para la tarea B es:

\begin{lstlisting}
L_B(theta) = L_B_orig(theta) + (lambda/2) * Sum_i F_i * (theta_i - theta*_A,i)^2
\end{lstlisting}

donde $\theta^*_A$ son los pesos óptimos para la tarea A, $F_i$ es la importancia del peso $i$, y
$\lambda$ controla el equilibrio entre plasticidad y estabilidad. \textbf{SI} (\textit{Synaptic
Intelligence}, Zenke et al., 2017)~\cite{ref5} estima las importancias de forma online durante el
entrenamiento, evitando el coste computacional del cálculo de la Fisher.

El principal inconveniente de estos métodos es que la penalización es solo una aproximación: no
garantizan la ausencia total de olvido, y la importancia calculada sobre una tarea puede no ser
compatible con las importancias calculadas sobre múltiples tareas sucesivas.

\subsection{Métodos basados en repetición (\textit{replay})}

Estos métodos mantienen una memoria de ejemplos pasados (reales o generados) que se mezclan con los
datos nuevos durante el re-entrenamiento. El \textbf{replay exacto} almacena directamente
subconjuntos de datos históricos (\textit{coreset methods}), con el coste de violar potencialmente
restricciones de privacidad o almacenamiento. El \textbf{replay generativo} (Shin et
al., 2017)~\cite{ref10} usa un modelo generativo (típicamente una GAN) entrenado sobre las tareas
pasadas para sintetizar ejemplos artificiales que se mezclan con los datos nuevos.

El replay generativo elimina la necesidad de almacenar datos reales, pero introduce un nuevo modelo
que debe entrenarse a su vez de forma continua, trasladando el problema del olvido al generador.

\subsection{Métodos basados en arquitectura}

Estos métodos modifican la estructura de la red para aislar el conocimiento de cada tarea. Las
\textbf{Redes Neuronales Progresivas} (PNN, Rusu et al., 2016)~\cite{ref6} añaden una columna nueva
a la red para cada tarea nueva, con conexiones laterales a las columnas anteriores, pero congelando
los pesos previos. Esto garantiza que las tareas antiguas no se olvidan, pero el tamaño del modelo
crece linealmente con el número de tareas.

Los métodos basados en máscaras binarias (\textit{PackNet}, Mallya y Lazebnik, 2018~\cite{ref11};
\textit{HAT}, Serrà et al., 2018~\cite{ref12}) reservan subconjuntos de pesos para cada tarea, pero
requieren conocer el número de tareas de antemano y tienen una capacidad finita.

\section{Linealidad local de redes ReLU}

Una propiedad matemática fundamental de las redes con activación ReLU es que definen funciones
\textbf{continuas y lineales a tramos}: para una entrada dada, la composición de transformaciones
afines y funciones ReLU resulta en una transformación lineal afín en la región convexa del espacio
de entrada que contiene a esa muestra~\cite{ref3}. Esta región queda determinada por el
\textbf{patrón de activación} binario: el vector que indica qué neuronas tienen salida positiva y
cuáles cero.

Esta propiedad ha sido estudiada desde la perspectiva de la capacidad representacional (las redes
ReLU profundas pueden generar un número de regiones lineales exponencial con la
profundidad~\cite{ref3}) y más recientemente desde la perspectiva del análisis de
robustez~\cite{ref13}. Qsimov explota esta propiedad de forma constructiva: en lugar de analizar
las regiones lineales de la red completa, extrae y enumera los \textbf{caminos de conexión activos}
entre capas, que son la representación dual de las regiones lineales cuando se consideran las
conexiones capa a capa.

\section{Sistemas lineales para re-entrenamiento}

La idea de reformular el entrenamiento de redes neuronales como un problema de mínimos cuadrados es
clásica en el contexto de las capas de salida. La técnica \textbf{LWTA} (\textit{Local Winner Takes
All}, Srivastava et al., 2013~\cite{ref14}) y los métodos de \textbf{extreme learning machines}
(ELM, Huang et al., 2006~\cite{ref15}) fijan los pesos de las capas ocultas aleatoriamente y
entrenan solo la última capa con mínimos cuadrados. Qsimov se distingue de estos enfoques en que
los pesos del modelo base no son aleatorios sino pre-entrenados, y en que la representación lineal
no se aplica solo a la última capa sino a una secuencia de capas cuyo punto de inicio es
configurable.

\section{Resumen y posicionamiento de Qsimov}

La tabla siguiente sitúa Qsimov en relación con las principales aproximaciones al aprendizaje
continuo:

\begin{table}[htbp]
  \centering
  \caption{Comparativa de métodos de aprendizaje continuo.}
  \label{tab:comparativa}
  \resizebox{\textwidth}{!}{%
  \begin{tabular}{lcccc}
    \toprule
    \textbf{Método} & \textbf{Sin datos previos} & \textbf{Sin arq. creciente} & \textbf{Garantía no-olvido} & \textbf{Velocidad} \\
    \midrule
    Fine-tuning estándar            & \checkmark & \checkmark & $\times$              & Alta \\
    EWC~\cite{ref4}                 & \checkmark & \checkmark & Parcial               & Media \\
    Replay exacto                   & $\times$   & \checkmark & Parcial               & Baja \\
    Replay generativo~\cite{ref10}  & \checkmark & $\times$   & Parcial               & Baja \\
    PNN~\cite{ref6}                 & \checkmark & $\times$   & \checkmark            & Media \\
    \textbf{Qsimov (sist. lineal)}  & \checkmark & \checkmark & \checkmark (datos vis.) & Muy alta \\
    \bottomrule
  \end{tabular}}%
\end{table}

Qsimov es el único método en esta comparativa que combina ausencia de necesidad de datos históricos,
tamaño de modelo constante, garantía exacta de no-olvido para los datos incluidos, y actualización
de coste prácticamente constante independientemente del histórico acumulado.

El hueco que ocupa Qsimov en el panorama actual es, por tanto, la intersección de tres propiedades
que hasta ahora no se daban simultáneamente en ningún método revisado: no necesitar acceso a datos
pasados, no aumentar la arquitectura y ofrecer garantía formal de no-olvido. El presente trabajo
evalúa empíricamente si estas propiedades teóricas se traducen en ventajas medibles en escenarios de
referencia estándar.


\chapter{Especificación del Sistema}
% Capítulo 3: Especificación del Sistema

\section{Definición de requisitos}

\subsection{Requisitos funcionales}

\paragraph{RF-01} El sistema debe permitir seleccionar un punto de corte en la red base (capa
\texttt{initial\_layer}) a partir del cual se aplica el algoritmo Qsimov.

\paragraph{RF-02} El sistema debe enumerar todos los caminos activos entre la capa
\texttt{initial\_layer} y la salida de la red, para capas de tipo Dense, Conv1D/2D/3D y
MaxPooling1D/2D/3D con activación ReLU.

\paragraph{RF-03} Dado un conjunto de muestras de entrenamiento, el sistema debe generar el sistema
de ecuaciones lineales $Ax = b$ donde cada fila representa la contribución de una muestra a los
pesos de los caminos.

\paragraph{RF-04} El sistema debe resolver el sistema lineal mediante back-substitution (solución
exacta) o mínimos cuadrados (solución aproximada).

\paragraph{RF-05} El sistema debe permitir acumular ecuaciones de distintas rondas de entrenamiento
sin resetear el sistema, implementando así el mecanismo de no-olvido.

\paragraph{RF-06} El sistema debe implementar una variante basada en descenso de gradiente sobre
una capa densa restringida a los caminos activos.

\paragraph{RF-07} El sistema debe ser compatible con los frameworks TensorFlow/Keras y PyTorch.

\paragraph{RF-08} El sistema debe permitir guardar y cargar modelos entrenados en un formato propio
(directorios \texttt{.qsi}).

\subsection{Requisitos no funcionales}

\paragraph{RNF-01 (Rendimiento)} Las operaciones de extracción de coeficientes y generación de
caminos deben estar optimizadas con extensiones Cython compiladas.

\paragraph{RNF-02 (Compatibilidad)} La librería debe funcionar con Python 3.8 y las versiones
TensorFlow 2.11 y PyTorch 2.0.

\paragraph{RNF-03 (Modularidad)} Las implementaciones para Keras y PyTorch deben compartir clases
base abstractas que definan la interfaz pública.

\paragraph{RNF-04 (Testeabilidad)} Todas las clases principales deben tener cobertura de tests
unitarios e integración ejecutables con \texttt{pytest}.

\section{Especificación del sistema}

El sistema se organiza en torno a tres componentes principales, cuya relación se describe a
continuación:

\begin{figure}[htbp]
  \centering
\begin{verbatim}
              +-------------------------------+
              |   Red base pre-entrenada      |
              | (Keras Sequential / PyTorch)  |
              +---------------+---------------+
                              |
              +---------------v---------------+
              |         PathSelector           |
              |  - left_model_  (hasta i_layer)|
              |  - right_model_ (desde i_layer)|
              |  - all_paths_                  |
              |  - samples_to_coefficients()   |
              +---------------+---------------+
                              |
          +-------------------+-------------------+
          |                                       |
+---------v-----------+               +-----------v-----------+
|  QsimovLinearSystem |               |    QsimovGradient     |
|  - fit() (acumul.)  |               |  - fit() (por epochs) |
|  - predict()        |               |  - predict()          |
|  - reset_equations()|               |  - model_ (CustomDens)|
+---------------------+               +-----------------------+
\end{verbatim}
  \caption{Diagrama de componentes del sistema Qsimov.}
  \label{fig:componentes}
\end{figure}

\section{Planificación y estimación de costes}

\subsection{Planificación}

El desarrollo del proyecto se ha estructurado en cinco fases principales:

\begin{table}[htbp]
  \centering
  \caption{Planificación del proyecto.}
  \label{tab:planificacion}
  \begin{tabular}{cp{8.5cm}c}
    \toprule
    \textbf{Fase} & \textbf{Descripción} & \textbf{Duración estimada} \\
    \midrule
    F1 & Revisión bibliográfica y estudio del estado del arte                       & 3 semanas \\
    F2 & Comprensión de la librería Qsimov y configuración del entorno experimental & 3 semanas \\
    F3 & Diseño de los experimentos y selección de configuraciones                  & 2 semanas \\
    F4 & Ejecución de experimentos y análisis de resultados                         & 8 semanas \\
    F5 & Redacción de la memoria y preparación de la defensa                        & 3 semanas \\
    \midrule
    \textbf{Total} & & \textbf{19 semanas} \\
    \bottomrule
  \end{tabular}
\end{table}

\subsection{Estimación de costes}

\textbf{Coste de personal} (estimación a precio de junior en prácticas, 15\,€/h; tarifa orientativa
para el sector TIC en España, acorde al convenio colectivo de empresas de ingeniería y consultoría
para titulados recientes):
\begin{itemize}
  \item 19 semanas $\times$ 25 h/semana = 475\,h
  \item Coste: 475\,h $\times$ 15\,€/h = \textbf{7.125\,€}
\end{itemize}

\textbf{Coste de hardware}:
\begin{itemize}
  \item Uso de GPU (NVIDIA, estimación de servidor cloud): $\sim$80\,h $\times$ 0.5\,€/h =
        \textbf{40\,€}
\end{itemize}

\textbf{Software}: Todas las herramientas utilizadas (Python, TensorFlow, PyTorch, Cython) son de
código abierto y gratuitas. Coste: \textbf{0\,€}

\textbf{Coste total estimado: 7.165\,€}


\chapter{Desarrollo del Proyecto}
% Capítulo 4: Desarrollo del Proyecto

\section{Análisis del sistema}

\subsection{Formalización matemática de los caminos}

Sea una red neuronal feed-forward con activación ReLU. Para una entrada $x \in \mathbb{R}^n$, la
activación de cada neurona en cada capa es no-negativa (por definición de ReLU). Se dice que la
neurona $j$ de la capa $l$ está \textbf{activa} para $x$ si su pre-activación es estrictamente
positiva, e \textbf{inactiva} si es cero.

Un \textbf{camino} (\textit{path}) $p = (j_0, j_1, \ldots, j_L)$ es una secuencia de índices de
neuronas, una por capa, tal que la neurona $j_l$ de la capa $l$ está conectada con la neurona
$j_{l+1}$ de la capa $l+1$. El \textbf{valor del coeficiente} de un camino para una entrada $x$ es
el producto de:
\begin{itemize}
  \item La activación de la neurona de entrada del camino (neurona $j_0$), y
  \item El indicador de actividad de cada neurona intermedia (1 si activa, 0 si no).
\end{itemize}

Dado el conjunto de todos los caminos posibles entre la capa \texttt{initial\_layer} y la salida, y
dado un conjunto de muestras $\{(x_i, y_i)\}$, el sistema lineal $Ax = b$ se construye de la
siguiente forma:
\begin{itemize}
  \item Cada fila de $A$ corresponde a una muestra $x_i$ y contiene los coeficientes de los
        caminos activos para esa muestra.
  \item $b$ contiene las salidas deseadas $y_i$.
  \item $x$ son los pesos de los caminos (incógnitas del problema).
\end{itemize}

La solución $x^*$ del sistema proporciona los pesos óptimos que minimizan el error cuadrático entre
la salida de la red reconstruida y las salidas deseadas.

\subsection{Casos de uso principales}

\paragraph{CU-01: Entrenar Qsimov (primera ronda)}
\begin{itemize}
  \item Actor: usuario de la librería
  \item Precondición: modelo Keras/PyTorch pre-entrenado disponible
  \item Flujo: (1) crear \texttt{PathSelector} con \texttt{initial\_layer}, (2) crear
        \texttt{QsimovLinearSystem}, (3) llamar a \texttt{fit(X\_train, Y\_train)}
  \item Postcondición: \texttt{solutions\_} contiene los pesos de los caminos resolventes
\end{itemize}

\paragraph{CU-02: Actualizar el modelo sin olvido (ronda N)}
\begin{itemize}
  \item Actor: usuario
  \item Precondición: \texttt{QsimovLinearSystem} entrenado en rondas anteriores
  \item Flujo: llamar a \texttt{fit(X\_nueva, Y\_nueva)}, \textit{sin} llamar
        previamente a \texttt{reset\_equations()}
  \item Postcondición: \texttt{solutions\_} incorpora todas las muestras de todas las rondas
\end{itemize}

\paragraph{CU-03: Predecir con el modelo actualizado}
\begin{itemize}
  \item Actor: usuario
  \item Precondición: \texttt{QsimovLinearSystem} con \texttt{solutions\_} calculadas
  \item Flujo: llamar a \texttt{predict(X\_test)}
  \item Postcondición: se devuelven las predicciones del modelo Qsimov
\end{itemize}

\paragraph{CU-04: Guardar y recuperar el modelo}
\begin{itemize}
  \item Actor: usuario
  \item Guardar: \texttt{qls.save("modelo.qsi")} \\
        Recuperar: \texttt{qls = KerasQsimov\-Linear\-System.load("modelo.qsi")}
\end{itemize}

\subsection{Modelado de datos}

Los datos que gestiona el sistema son:
\begin{itemize}
  \item \texttt{all\_paths\_}: lista de arrays numpy de forma \texttt{(n\_paths\_layer, path\_length)},
        los índices que forman cada camino entre capas.
  \item \texttt{output\_masks\_}: array booleano
        \texttt{(n\_outputs, n\_paths\_total)};
        indica qué caminos llegan a cada neurona de salida.
  \item \texttt{equations\_}: lista de sistemas \texttt{[A, b]} por output, donde \texttt{A} tiene
        forma \texttt{(n\_equations, n\_paths\_output)}.
  \item \texttt{solutions\_}: lista de arrays numpy de forma \texttt{(n\_paths\_output,)},
        pesos de los caminos resueltos.
\end{itemize}

\section{Diseño del sistema}

\subsection{Diseño de interfaces de usuario}

La API pública de la librería está diseñada para ser coherente con las convenciones de scikit-learn
(\texttt{fit} / \texttt{predict}) y extensible a nuevos frameworks. La interfaz mínima para el
usuario es:

\begin{lstlisting}[language=Python, caption={API mínima -- variante Keras.}]
# Keras
from qsimov.keras_path_selector import KerasPathSelector
from qsimov.keras_qsimov_linear_system import (
    KerasQsimovLinearSystem)

path_selector = KerasPathSelector(model, initial_layer=-1)
qls = KerasQsimovLinearSystem(
    path_selector, solver="back_substitution")
qls.fit(x_train, y_train)
y_pred = qls.predict(x_test)
\end{lstlisting}

\begin{lstlisting}[language=Python, caption={API mínima -- variante PyTorch.}]
# PyTorch
from qsimov.pytorch_path_selector import PytorchPathSelector
from qsimov.pytorch_qsimov_gradient import PytorchQsimovGradient

path_selector = PytorchPathSelector(
    model, input_shape=(1, 28, 28), initial_layer=-4)
qg = PytorchQsimovGradient(path_selector)
qg.fit(x_train, y_train, epochs=5,
       optimizer=lambda p: torch.optim.Adam(p, lr=1e-3))
\end{lstlisting}

\subsection{Diseño de procesos: extracción de caminos}

El proceso de extracción de caminos para una capa densa con matriz de pesos
$W \in \mathbb{R}^{m \times n}$ funciona de la siguiente forma:

\begin{enumerate}
  \item Para cada neurona de entrada $i$ (incluyendo el bias como neurona $0$): se recorren todas
        las neuronas de salida $j$ con $W[i, j] \neq 0$. Cada par $(i, j)$ constituye un camino
        elemental de longitud 2.
  \item Los caminos de capas sucesivas se combinan mediante \path{combine_paths_left_right}: se
        hace un join de los caminos izquierdos y derechos donde el último índice del camino
        izquierdo coincide con el primero del camino derecho. Los caminos de bias (con primer índice
        0) siempre se incluyen.
  \item El proceso se repite para todas las capas del \texttt{right\_model\_}, acumulando los
        caminos mediante combinación.
\end{enumerate}

Para capas convolucionales, cada posición del mapa de características y cada canal de salida genera
sus propios caminos, aplicando las transformaciones de stride y padding correspondientes.

\subsection{Diseño de la base de datos (persistencia)}

El formato de persistencia \texttt{.qsi} es un directorio con la siguiente estructura:

\begin{lstlisting}[caption={Estructura del directorio de persistencia \texttt{.qsi}.}]
modelo.qsi/
|-- py_objects.pkl           <- atributos escalares y listas (pickle)
|-- numpy_variables.npz      <- arrays numpy comprimidos (np.savez_compressed)
|-- left_model.h5 / .pt      <- left_model_ del PathSelector
|-- right_model.h5 / .pt     <- right_model_ del PathSelector
`-- path_selector.qsi/       <- PathSelector anidado
    |-- py_objects.pkl
    `-- numpy_variables.npz
\end{lstlisting}

Esta estructura separa los arrays numpy de gran tamaño (como \texttt{all\_paths\_} o
\texttt{equations\_}) del resto de atributos de Python, permitiendo una carga eficiente sin
deserializar objetos que no sean necesarios.

\section{Arquitectura interna de Qsimov}

\subsection{Módulo \texttt{paths/}: Generación de caminos (Python + Cython)}

Las operaciones de mayor coste computacional en la generación de caminos están implementadas en
Cython (\texttt{c\_paths.pyx}, \texttt{c\_combine.pyx}, \texttt{c\_dense.pyx},
\texttt{c\_conv.pyx}, \texttt{c\_maxpooling.pyx}). Cython se eligió sobre alternativas como NumPy
vectorizado o C puro porque las operaciones implican bucles anidados sobre arrays de longitud
variable con lógica de control compleja, donde la vectorización NumPy no es directamente aplicable.

La función central \texttt{c\_non\_zero\_input\_select\_paths} filtra, para cada muestra, los
caminos cuya neurona de entrada tiene activación no-nula. Su signatura en Cython es:

\begin{lstlisting}[language=C, caption={Signatura de \texttt{c\_non\_zero\_input\_select\_paths}.}]
def c_non_zero_input_select_paths(
    float[:] flat_inputs_with_bias,   # activaciones de entrada, con bias prepend
    long[:, :] all_paths,             # todos los caminos posibles
) -> tuple[np.ndarray, np.ndarray]:  # (selected_paths, select_mask)
\end{lstlisting}

\subsection{Módulo \texttt{linalg.py}: Resolución del sistema}

El sistema lineal se resuelve siguiendo estos pasos:

\begin{enumerate}
  \item \textbf{Transformación QR} (\texttt{r\_transform}): Se aplica la factorización QR con
        \texttt{numpy.linalg.qr(mode="r")} para obtener la matriz triangular superior $R$. Esto
        reduce el sistema acumulado, que puede tener miles de ecuaciones, a un sistema cuadrado
        de tamaño $n_{paths} \times n_{paths}$, manteniendo solo la información relevante.
  \item \textbf{Back-substitution} (\texttt{c\_back\_substitution}): Se resuelve
        $R \cdot x = Q^T \cdot b$ hacia atrás. Se aplican cutoffs absoluto y relativo para tratar
        los coeficientes casi-nulos como cero exacto, lo que mejora la estabilidad numérica.
  \item \textbf{Shrinkage incremental}: El parámetro \texttt{qr\_shrinkage\_factor} controla cada
        cuántas filas nuevas se aplica la transformación QR para comprimir el sistema en memoria.
\end{enumerate}

La variante de mínimos cuadrados (\texttt{solver="lstsq"}) delega directamente en
\texttt{numpy.linalg.lstsq}, útil cuando el sistema está sub-determinado (más incógnitas que
ecuaciones).

\subsection{Diferencias de implementación entre Keras y PyTorch}

Las dos implementaciones comparten las clases base abstractas \texttt{PathSelector} y
\texttt{QsimovLinearSystem}, pero difieren en los siguientes puntos:

\begin{table}[htbp]
  \centering
  \caption{Diferencias entre las variantes Keras y PyTorch de Qsimov.}
  \label{tab:tabla1}
  \begin{tabular}{p{5cm}p{4.5cm}p{4.5cm}}
    \toprule
    \textbf{Aspecto} & \textbf{Keras} & \textbf{PyTorch} \\
    \midrule
    Inferencia del \texttt{input\_shape}  & Automática (el modelo lo infiere) & Manual (parámetro requerido) \\
    Formato de pesos Dense                & $W \in \mathbb{R}^{in \times out}$ & $W \in \mathbb{R}^{out \times in}$ (se transpone) \\
    Formato de datos conv                 & \texttt{channels\_last}            & \texttt{channels\_first} \\
    Iterador de batches                   & \path{as_tensorflow_dataset()}  & \path{as_pytorch_dataloader()} \\
    Dispositivo del right model           & Siempre CPU (por diseño)           & Siempre CPU (por diseño) \\
    Guardado del modelo                   & Formato \texttt{.h5}               & Formato \texttt{.pt} \\
    Indexación de \texttt{initial\_layer} & Solo capas Dense/Conv (ReLU y Dropout no cuentan) & Todas las capas, incluidas ReLU y Dropout \\
    \bottomrule
  \end{tabular}
\end{table}


\chapter{Pruebas y Resultados}
% Capítulo 5: Pruebas y Resultados

\section{Descripción de experimentos}

Para la validación de Qsimov se han diseñado cinco experimentos, ordenados de menor a mayor
complejidad del dataset y del escenario evaluado.

\subsection{Configuración experimental común}

Todos los experimentos se ejecutan bajo las siguientes condiciones:

Todos los experimentos se ejecutaron en la GPU del servidor
\texttt{fermidi1-informat-uv-es} de la Universitat de València (NVIDIA GPU, Python 3.10). El
servidor dispone de acceso exclusivo a la GPU durante los experimentos, por lo que los tiempos de
entrenamiento no están distorsionados por la secuencialidad impuesta por el bajo número de cores de
CPU.

Se utilizó PyTorch 2.x como framework en todos los experimentos reportados. La librería Qsimov
ofrece implementaciones equivalentes en TensorFlow/Keras, pero todos los experimentos presentados en
este documento se han reejecutado con el backend PyTorch sobre GPU para garantizar comparativas de
tiempo homogéneas entre métodos. El entorno software concreto fue:
\begin{itemize}
  \item Python 3.10
  \item PyTorch 2.x con CUDA
  \item Cython 3.0 (extensiones compiladas con \texttt{python setup.py build\_ext --inplace})
  \item MLflow para tracking de experimentos
\end{itemize}

Las semillas aleatorias están fijadas en cada experimento. El tracking de experimentos se realiza
con MLflow~\cite{ref16}. El Anexo~C detalla los comandos exactos para reproducir cada experimento.

El modelo base se pre-entrena con el método estándar (Adam, cross-entropy) antes de aplicar Qsimov.

\subsection{Experimento 1: Velocidad y precisión (MNIST)}

\paragraph{Configuración} Dataset MNIST (60.000 muestras de entrenamiento, 10.000 de test; imágenes
de dígitos manuscritos $28 \times 28$ píxeles~\cite{ref17}). Red base: CNN con dos bloques
Conv2D+MaxPooling seguidos de tres capas Dense (18.474 parámetros totales). Punto de corte:
\texttt{initial\_layer = -2} (la última capa \texttt{Linear(16} $\rightarrow$ \texttt{10)} en PyTorch, ya que
ReLU y Dropout cuentan como capas separadas). Número de caminos generados: \textbf{17} (16 entradas
de \texttt{Linear(16} $\rightarrow$ \texttt{10)} más el nodo de bias).

\begin{table}[htbp]
  \centering
  \caption{Métodos comparados en el experimento de velocidad y precisión (MNIST).}
  \label{tab:metodos_mnist}
  \begin{tabular}{lp{9cm}}
    \toprule
    \textbf{Método} & \textbf{Descripción} \\
    \midrule
    \texttt{standard} (half) & Re-entrenamiento estándar PyTorch con la mitad del dataset (30.000 muestras), 10 epochs \\
    \texttt{standard} (full) & Re-entrenamiento estándar PyTorch con el dataset completo (60.000 muestras), 10 epochs \\
    \texttt{qsimov\_gradient} & QsimovGradient sobre la mitad del dataset (PathSelector), ajustado en el dataset completo, 10 epochs \\
    \texttt{qsimov\_linear}   & QsimovLinearSystem sobre la mitad del dataset (PathSelector), ajustado en el dataset completo, una sola pasada \\
    \bottomrule
  \end{tabular}
\end{table}

Se comparan dos funciones de pérdida: entropía cruzada categórica y MSE. El \texttt{qsimov\_linear}
solo se evalúa con MSE, puesto que el sistema lineal requiere que la función objetivo sea cuadrática.

\begin{table}[htbp]
  \centering
  \caption{Resultados del experimento de velocidad y precisión sobre MNIST (GPU, PyTorch, MSE).}
  \label{tab:tabla2}
  \begin{tabular}{lcc}
    \toprule
    \textbf{Método} & \textbf{Precisión test (\%)} & \textbf{Tiempo (s)} \\
    \midrule
    \texttt{standard} (half)              & 97.42 & 72.4 \\
    \texttt{standard} (full)              & 97.32 & 144.9 \\
    \texttt{qsimov\_gradient} (half$\rightarrow$full) & 97.40 & 90.5 \\
    \texttt{qsimov\_linear} (half$\rightarrow$full)   & 97.40 & 6.6 \\
    \bottomrule
  \end{tabular}
\end{table}

El sistema lineal de Qsimov alcanza una precisión equivalente al re-entrenamiento estándar
entrenado con el dataset completo, pero en un tiempo \textbf{22 veces menor} (6.6\,s frente a
144.9\,s). Esta ventaja temporal se debe a que \texttt{QsimovLinearSystem} no realiza descenso de
gradiente iterativo: en su lugar, resuelve el sistema de ecuaciones $Ax = b$ en una única pasada
sobre los datos mediante factorización QR y back-substitution. La diferencia de precisión entre
todos los métodos es inferior a 0.1 puntos porcentuales, lo que indica que la reducción del
re-entrenamiento a los caminos activos de la última capa no implica pérdida significativa de
capacidad representacional.

\subsection{Experimento 2: Robustez al learning rate (MNIST)}

Los resultados anteriores confirmaron que el sistema lineal iguala la precisión del
re-entrenamiento estándar. El siguiente experimento se centra en el módulo de gradiente y en una
debilidad práctica frecuente del re-entrenamiento estándar: la sensibilidad al learning rate, cuya
elección incorrecta puede degradar drásticamente el rendimiento.

\paragraph{Configuración} Mismo dataset y red que el experimento anterior (PyTorch, GPU). Se evalúa
el rendimiento de \texttt{qsimov\_gradient} y del re-entrenamiento estándar PyTorch en un rango de
6 learning rates logarítmicamente espaciados ($6.25 \times 10^{-5}$, $2.5 \times 10^{-4}$,
$1 \times 10^{-3}$, $4 \times 10^{-3}$, $1.6 \times 10^{-2}$, $6.4 \times 10^{-2}$). Ambos
métodos entrenan durante 10 epochs.

\begin{table}[htbp]
  \centering
  \caption{Precisión de \texttt{qsimov\_gradient} y re-entrenamiento estándar en función del
           learning rate (GPU, PyTorch).}
  \label{tab:tabla3}
  \begin{tabular}{ccc}
    \toprule
    \textbf{Learning rate} & \textbf{\texttt{qsimov\_gradient} (\%)} & \textbf{\texttt{standard} PyTorch (\%)} \\
    \midrule
    $6.25 \times 10^{-5}$ & 94.3 & 87.3 \\
    $2.5  \times 10^{-4}$ & 97.5 & 96.4 \\
    $1    \times 10^{-3}$ & 97.9 & 97.7 \\
    $4    \times 10^{-3}$ & 98.1 & 98.2 \\
    $1.6  \times 10^{-2}$ & 98.2 & 97.0 \\
    $6.4  \times 10^{-2}$ & 97.8 & 44.4 \\
    \bottomrule
  \end{tabular}
\end{table}

La red estándar degrada significativamente con el learning rate más alto (44.4\,\%), mientras que
\texttt{qsimov\_gradient} se mantiene en el rango 94--98\,\% en todo el espectro. Esta robustez se
explica porque la capa \texttt{CustomConnectedLinear} del módulo de gradiente tiene muchas menos
conexiones que la capa Dense original (solo los caminos activos), lo que reduce el espacio de
búsqueda y estabiliza el entrenamiento ante tasas de aprendizaje elevadas.

\subsection{Experimento 3: Olvido catastrófico (MNIST)}

La velocidad y la robustez al learning rate quedan establecidas en los experimentos anteriores. La
pregunta central del TFG es si Qsimov resuelve el problema del olvido catastrófico. Este
experimento lo verifica directamente: un modelo entrenado sobre los dígitos 0--4 se actualiza con
datos exclusivamente de los dígitos 5--9, sin acceso a los datos originales.

\paragraph{Configuración} MNIST dividido en dos fases: fase 1 con los dígitos 0--4 (30.596 train,
5.139 test), fase 2 con los dígitos 5--9 (29.404 train, 4.861 test). Ejecutado con PyTorch sobre
GPU. El modelo base se pre-entrena sobre la fase 1 durante 10 epochs. En la fase 2 no se permite
acceder a los datos de la fase 1 (excepto en el método \texttt{standard\_cumulative}, que sirve de
cota superior teórica).

\begin{table}[htbp]
  \centering
  \caption{Métodos comparados en el experimento de olvido catastrófico (MNIST).}
  \label{tab:metodos_forgetting_mnist}
  \begin{tabular}{lcp{7cm}}
    \toprule
    \textbf{Método} & \textbf{Acc. fase 1} & \textbf{Descripción} \\
    \midrule
    \texttt{base}              & ---        & Modelo original, sin actualización (referencia) \\
    \texttt{linear\_accum}     & No         & QsimovLinearSystem: fit fase 1 + fit fase 2 sin reset \\
    \texttt{linear\_new\_only} & No         & QsimovLinearSystem: solo datos de fase 2 (reset entre fases) \\
    \texttt{gradient\_new\_only} & No       & QsimovGradient: 5 epochs solo en datos de fase 2 \\
    \texttt{finetune\_new\_only} & No       & Fine-tuning estándar PyTorch: 5 epochs solo en datos de fase 2 \\
    \texttt{cumulative}        & \textbf{Sí} & Fine-tuning estándar con datos de fase 1 + fase 2 (acumulado) \\
    \bottomrule
  \end{tabular}
\end{table}

\begin{table}[htbp]
  \centering
  \caption{Resultados del experimento de olvido catastrófico en MNIST (GPU, PyTorch).}
  \label{tab:tabla4}
  \small
  \begin{tabular}{lcccc}
    \toprule
    \textbf{Método} & \textbf{Acc. 0--4 (\%)} & \textbf{Acc. 5--9 (\%)} & \textbf{Acc. global (\%)} & \textbf{Tiempo (s)} \\
    \midrule
    \texttt{base}              &  99.6        &  0.0         & 51.2        & 0.0 \\
    \texttt{linear\_accum}     & \textbf{94.3} & \textbf{49.1} & \textbf{72.3} & \textbf{1.2} \\
    \texttt{linear\_new\_only} &   0.0        & 68.5         & 33.3        & 1.1 \\
    \texttt{gradient\_new\_only} & 0.0        & 61.1         & 29.7        & 7.7 \\
    \texttt{finetune\_new\_only} & 0.0        & 96.1         & 46.7        & 6.4 \\
    \texttt{cumulative}        &  98.5        & 94.8         & 96.7        & 12.6 \\
    \bottomrule
  \end{tabular}
\end{table}

\texttt{linear\_accum} es el único método que mantiene precisión en las clases antiguas (94.3\,\%)
mientras aprende las nuevas (49.1\,\%), y lo hace en \textbf{1.2\,s}, diez veces más rápido que
el re-entrenamiento acumulado y sin acceder a los datos de la fase 1. Los métodos
\texttt{linear\_new\_only}, \texttt{gradient\_new\_only} y \texttt{finetune\_new\_only} exhiben
olvido catastrófico total en las clases 0--4 (0\,\%). El re-entrenamiento acumulado
(\texttt{cumulative}) consigue el mejor rendimiento global (96.7\,\%) pero requiere acceso a los
datos históricos y tarda 12.6\,s. La precisión del \texttt{linear\_accum} en clases nuevas (49.1\,\%)
es inferior a la del fine-tuning sobre nuevas (96.1\,\%): este compromiso es el precio de no olvidar,
ya que el sistema lineal debe satisfacer simultáneamente las ecuaciones de las 30.596 muestras de la
fase 1 y las 29.404 de la fase 2.

\subsection{Experimento 4: CIFAR-10. Olvido catastrófico y datos escasos}

\paragraph{Configuración} Dataset CIFAR-10 (50.000 muestras de entrenamiento, 10.000 de test;
imágenes en color $32 \times 32$ píxeles~\cite{ref18}). Red base: LeNet adaptado con dos bloques
Conv2D+MaxPooling y tres capas Dense. Ejecutado con PyTorch sobre GPU. Punto de corte:
\texttt{initial\_layer = -1}. Número de caminos generados: \textbf{33}.

\textbf{4a.~Olvido catastrófico}: Modelo pre-entrenado en clases 0--4 (fase 1). Actualización con
clases 5--9 (fase 2) sin acceso a datos históricos. 20 epochs en fase 1, 5 epochs en fase 2.

\begin{table}[htbp]
  \centering
  \caption{Resultados del experimento de olvido catastrófico en CIFAR-10 (GPU, PyTorch).}
  \label{tab:tabla5}
  \resizebox{\textwidth}{!}{%
  \begin{tabular}{lcccc}
    \toprule
    \textbf{Método} & \textbf{Acc. antiguas (\%)} & \textbf{Acc. nuevas (\%)} & \textbf{Acc. global (\%)} & \textbf{Tiempo (s)} \\
    \midrule
    \texttt{base} (sin actualizar)   & 78.6         & 0.0           & 39.3        & 0.0 \\
    \texttt{linear\_accum}           & \textbf{73.4} & 8.7          & 41.1        & \textbf{1.3} \\
    \texttt{linear\_new\_only}       & 0.0          & 52.8          & 26.4        & 1.3 \\
    \texttt{gradient\_new\_only}     & 0.0          & 52.4          & 26.2        & 7.7 \\
    \texttt{finetune\_new\_only}     & 0.0          & \textbf{85.2} & 42.6        & 5.7 \\
    \texttt{cumulative} (acumulado)  & 62.7         & 72.9          & \textbf{67.8} & 11.6 \\
    \bottomrule
  \end{tabular}}%
\end{table}

\texttt{finetune\_new\_only} exhibe olvido catastrófico total: la precisión en las clases antiguas
cae de 78.6\,\% a 0.0\,\%. \texttt{linear\_accum} es el único método que preserva conocimiento de
las clases antiguas (73.4\,\%) sin acceso a datos históricos, y lo hace en solo 1.3\,s. La
precisión en clases nuevas es baja (8.7\,\%) porque con solo 33 caminos el espacio representacional
del sistema lineal es muy reducido: resolver simultáneamente las ecuaciones de la fase 1 y de la
fase 2 deja pocos grados de libertad para las clases nuevas. Los métodos
\texttt{linear\_new\_only} y \texttt{gradient\_new\_only} obtienen $\sim$52\,\% en clases nuevas
pero olvidan completamente las antiguas, comportándose como el fine-tuning estándar.

\textbf{4b.~Generalización con datos escasos}: Mismo modelo LeNet entrenado con distintos
subconjuntos del dataset completo (100, 1.000 y 5.000 muestras), comparando gradiente estándar
versus \texttt{qsimov\_gradient}. 10 epochs en todos los casos.

\begin{table}[htbp]
  \centering
  \caption{Precisión en test según el tamaño del split de entrenamiento (GPU, PyTorch, 10 epochs).}
  \label{tab:tabla5b}
  \small
  \begin{tabular}{p{4cm}cc}
    \toprule
    \textbf{Tamaño del split} & \textbf{Std. gradient (\%)} & \textbf{Qsimov gradient (\%)} \\
    \midrule
    100 muestras              & 10.7          & \textbf{20.0} \\
    1.000 muestras            & 25.7          & \textbf{29.8} \\
    5.000 muestras            & 36.8          & 36.0 \\
    Dataset completo (50.000) & \textbf{54.9} & --- \\
    \bottomrule
  \end{tabular}
\end{table}

Con conjuntos de datos muy pequeños (100 muestras), \texttt{qsimov\_gradient} supera al gradiente
estándar en $\sim$9 puntos porcentuales (20.0\,\% vs.\ 10.7\,\%). Esta ventaja se reduce con más
datos y prácticamente desaparece en torno a las 5.000 muestras (36.8\,\% vs.\ 36.0\,\%), donde la
mayor capacidad del gradiente estándar sobre toda la red compensa la regularización implícita de
Qsimov. El resultado confirma que la restricción estructural de Qsimov es una regularización
implícita efectiva en el régimen de escasez de datos.

\subsection{Experimento 5: Aprendizaje continual en ImageNet (100 clases)}

\paragraph{Configuración} Subconjunto de 100 clases de ImageNet (130.000 imágenes de entrenamiento,
$\sim$13.000 de test). Red base: VGG16~\cite{ref19} con la última capa adaptada a 100 clases
(pre-entrenado sobre estas 100 clases desde cero). El escenario de aprendizaje continual divide los
datos en 4 rondas de 25 clases nuevas cada una. Los cuatro métodos comparados fueron
\texttt{qsimov\_linear\_accum}, \texttt{qsimov\_gradient}, \texttt{standard\_finetune} y
\texttt{standard\_cumulative}.

\begin{table}[htbp]
  \centering
  \caption{Resultados globales del experimento de aprendizaje continual sobre ImageNet (100 clases).}
  \label{tab:tabla6}
  \begin{tabular}{lcc}
    \toprule
    \textbf{Método} & \textbf{Acc. global (\%)} & \textbf{Tiempo total (s)} \\
    \midrule
    \texttt{qsimov\_linear\_accum}       & \textbf{37.0} & \textbf{113} \\
    \texttt{standard\_cumulative} (acum.)& 23.0          & 655 \\
    \texttt{standard\_finetune}          & 20.8          & --- \\
    \texttt{qsimov\_gradient}            &  2.5          & --- \\
    \bottomrule
  \end{tabular}
\end{table}

Qsimov con acumulación de ecuaciones supera al re-entrenamiento acumulado en precisión global
(37.0\,\% vs.\ 23.0\,\%) y lo hace seis veces más rápido (113\,s vs.\ 655\,s). La preservación del
conocimiento entre rondas es notable: la precisión en las rondas anteriores se mantiene entre el
79\,\% y el 84\,\% a lo largo de todo el experimento.

El módulo de gradiente de Qsimov no converge en este escenario (2.5\,\% global): con 100 clases y
un learning rate reducido, el método iterativo requeriría muchas más epochs de las utilizadas. Esta
es una limitación identificada del componente de gradiente en problemas de alta complejidad.

\section{Resultados y discusión}

Los cinco experimentos permiten extraer las siguientes conclusiones cuantitativas:

\begin{enumerate}
  \item El sistema lineal de Qsimov elimina el olvido catastrófico en todos los escenarios
        evaluados. En MNIST, \texttt{linear\_accum} mantiene 94.3\,\% en clases antiguas mientras
        aprende un 49.1\,\% de las nuevas, en solo 1.2\,s. En CIFAR-10, preserva 73.4\,\% en
        clases antiguas frente al 0.0\,\% del fine-tuning estándar, en 1.3\,s. En ImageNet (100
        clases, 4 rondas), termina con 37.0\,\% de precisión global, superando al re-entrenamiento
        acumulado (23.0\,\%) y a una velocidad $6\times$ superior.

  \item La elección del \texttt{initial\_layer} condiciona la capacidad del sistema. En
        MNIST, \texttt{initial\_layer = -2} restringe el re-entrenamiento a la última capa
        \texttt{Linear(16} $\rightarrow$ \texttt{10)}, generando 17 caminos; pese a ser pocos, resultan
        suficientes para una tarea relativamente simple (dígitos $28 \times 28$) y el sistema
        alcanza un 49.1\,\% en clases nuevas. En CIFAR-10, \texttt{initial\_layer = -1} genera 33
        caminos para un problema más complejo (imágenes naturales en color), pero ese espacio
        representacional resulta insuficiente: resolver simultáneamente las ecuaciones de la fase 1
        y las de la fase 2 deja pocos grados de libertad para las clases nuevas (8.7\,\%). La
        idoneidad del punto de corte depende, por tanto, de la relación entre el número de caminos
        generados y la complejidad intrínseca del problema, lo que convierte al
        \texttt{initial\_layer} en el hiperparámetro más crítico del sistema.

  \item El sistema lineal es más rápido que el re-entrenamiento estándar tanto en escenarios de re-entrenamiento puro ($\times 22$ en MNIST respecto al
        entrenamiento estándar completo; $\times 9$ en CIFAR-10 respecto al re-entrenamiento
        acumulado) como en aprendizaje continual ($\times 10$ en MNIST respecto al re-entrenamiento
        acumulado; $\times 6$ en ImageNet).

  \item El módulo de gradiente es una regularización implícita efectiva en regímenes de
        escasez de datos: con 100 muestras en CIFAR-10, \texttt{qsimov\_gradient} supera al
        gradiente estándar en $\sim$9 puntos (20.0\,\% vs.\ 10.7\,\%). La ventaja desaparece con
        conjuntos grandes donde la mayor capacidad del gradiente estándar domina (5.000 muestras:
        36.8\,\% vs.\ 36.0\,\%).

  \item El módulo de gradiente es robusto al learning rate en problemas de baja
        complejidad (MNIST: 94--98\,\% en todo el rango de LRs evaluado) pero no escala bien a
        problemas de alta complejidad (ImageNet, 100 clases: 2.5\,\% global con
        LR$=1\times 10^{-5}$ y 10 epochs).
\end{enumerate}

\section{Evaluación presupuestaria}

\begin{table}[htbp]
  \centering
  \caption{Evaluación presupuestaria: estimado vs.\ real.}
  \label{tab:presupuesto}
  \begin{tabular}{lcc}
    \toprule
    \textbf{Concepto} & \textbf{Estimado} & \textbf{Real} \\
    \midrule
    Horas de desarrollo           & 475\,h    & $\sim$490\,h \\
    Coste de personal (15\,€/h)   & 7.125\,€  & 7.350\,€ \\
    Coste GPU (cloud)             & 40\,€     & $\sim$55\,€ \\
    \midrule
    \textbf{Total}                & \textbf{7.165\,€} & \textbf{$\sim$7.405\,€} \\
    \bottomrule
  \end{tabular}
\end{table}

La desviación respecto a la estimación es de un 3,4\,\%, dentro del margen habitual para proyectos
de esta naturaleza.


\chapter{Conclusiones y Trabajo Futuro}
% Capítulo 6: Conclusiones y Trabajo Futuro

\section{Conclusiones}

Este trabajo ha evaluado experimentalmente Qsimov, una librería Python para re-entrenamiento
incremental sin olvido catastrófico en redes neuronales con activación ReLU. La contribución central
de este TFG es la caracterización experimental del comportamiento de Qsimov sobre MNIST, CIFAR-10
e ImageNet (100 clases), analizando las condiciones bajo las cuales su mecanismo de acumulación de
ecuaciones lineales elimina el olvido catastrófico.

Los resultados confirman que:

\begin{itemize}
  \item El sistema lineal de Qsimov supera a todos los métodos comparados en el escenario de
        aprendizaje continual sobre ImageNet, con un 37.0\,\% de precisión global frente al 23.0\,\%
        del re-entrenamiento acumulado y el 20.8\,\% del fine-tuning estándar.
  \item El tiempo de actualización del sistema lineal permanece prácticamente constante
        ($\approx$5.8\,s por batch en el experimento de streaming de ImageNet),
        independientemente del número de clases ya aprendidas.
  \item La robustez al learning rate del módulo de gradiente es notable en problemas de baja
        complejidad, eliminando la necesidad de ajuste fino de este hiperparámetro.
\end{itemize}

El principal factor limitante identificado es el número de caminos generados por el PathSelector,
que depende directamente del punto de corte elegido y de la densidad de la red base. En redes con
pocas neuronas antes de la capa de salida (CIFAR-10 con 33 caminos), el espacio representacional
puede no ser suficiente para aprender nuevas clases sin comprometer la capacidad en las antiguas.

Los resultados obtenidos son representativos de redes feed-forward con activación ReLU que incluyen
capas densas en su parte final, que es la arquitectura para la que Qsimov está diseñado. Cabe
esperar un comportamiento análogo en otras redes que cumplan estas condiciones, dado que el
mecanismo de acumulación de ecuaciones es agnóstico al dominio del problema. La extrapolación a
arquitecturas con conexiones residuales (ResNet, EfficientNet) o a modelos basados en atención
(Transformers) requeriría extender el módulo de generación de caminos y no puede garantizarse sin
evaluación experimental adicional.

\section{Trabajo futuro}

Las siguientes líneas de mejora se identifican como prioritarias para el trabajo futuro:

\begin{enumerate}
  \item \textbf{Selección automática del \texttt{initial\_layer}}: Diseñar un criterio cuantitativo
        (por ejemplo, basado en el número de caminos o en la complejidad del problema) que guíe la
        elección del punto de corte de forma automática, reduciendo la necesidad de conocimiento
        experto sobre la arquitectura.

  \item \textbf{Soporte para arquitecturas con conexiones residuales} (ResNet, EfficientNet): Las
        conexiones \textit{skip} rompen la estructura secuencial de caminos asumida por la
        implementación actual. Extender el módulo \texttt{paths/} para manejar estas arquitecturas
        ampliaría considerablemente el ámbito de aplicación.

  \item \textbf{Métricas de evaluación más completas}: Los experimentos actuales se centran en
        precisión global y tiempo. Añadir métricas de uso de memoria, análisis de la distribución
        de pesos de caminos, y curvas de estabilidad-plasticidad permitiría una evaluación más
        rigurosa.

  \item \textbf{Convergencia del módulo de gradiente en alta complejidad}: Explorar estrategias de
        warm-start, normalización por lotes, o ajuste adaptativo del learning rate para hacer
        converger el módulo de gradiente en escenarios de 100 o más clases.

  \item \textbf{Benchmarks adicionales}: Evaluar Qsimov sobre benchmarks estándar de aprendizaje
        continual como Split-CIFAR-100, Permuted-MNIST, o CORe50, para comparar cuantitativamente
        con métodos como EWC o iCaRL en condiciones homologables.
\end{enumerate}


\pagestyle{appendix}

\appendix
\chapter{Anexos}
% Apéndices

\section{Anexo A: Instalación y configuración del entorno}

\begin{lstlisting}[language=bash, caption={Instalación del entorno de desarrollo.}]
# Clonar el repositorio
git clone <repo_url>
cd qsimov

# Crear entorno virtual con Python 3.8
python3.8 -m venv venv
source venv/bin/activate

# Instalar dependencias
pip install -r requirements.txt

# Compilar extensiones Cython (obligatorio antes del primer uso)
python setup.py build_ext --inplace
\end{lstlisting}

\section{Anexo B: Ejemplo de uso completo (Keras, aprendizaje continual)}

\begin{lstlisting}[language=Python, caption={Ejemplo completo de aprendizaje continual con Keras.}]
import numpy as np
import tensorflow.keras as kr
from qsimov.keras_path_selector import KerasPathSelector
from qsimov.keras_qsimov_linear_system import KerasQsimovLinearSystem

# 1. Pre-entrenar modelo con datos de la fase 1
model = kr.Sequential([
    kr.layers.Conv2D(32, 3, activation='relu', input_shape=(28, 28, 1)),
    kr.layers.MaxPooling2D(),
    kr.layers.Flatten(),
    kr.layers.Dense(16, activation='relu'),
    kr.layers.Dense(10, activation='softmax'),
])
model.compile(optimizer='adam',
              loss='sparse_categorical_crossentropy',
              metrics=['accuracy'])
model.fit(x_phase1, y_phase1, epochs=10)

# 2. Crear PathSelector sobre el modelo entrenado
path_selector = KerasPathSelector(
    neural_network=model,
    initial_layer=-2,   # dos ultimas capas Dense
    verbose=1
)

# 3. Entrenar QsimovLinearSystem con fase 1
qls = KerasQsimovLinearSystem(
    path_selector=path_selector,
    solver="back_substitution",
    absolute_cutoff=1e-2,
    relative_cutoff=1e6,
)
qls.fit(x_phase1, y_phase1, batch_size=256)

# 4. Actualizar con fase 2 sin olvido (NO llamar reset_equations())
qls.fit(x_phase2, y_phase2, batch_size=256)

# 5. Evaluar en ambas fases
acc_old = np.mean(qls.predict(x_test_phase1).argmax(1) == y_test_phase1)
acc_new = np.mean(qls.predict(x_test_phase2).argmax(1) == y_test_phase2)
print(f"Precision clases antiguas: {acc_old:.3f}")
print(f"Precision clases nuevas:   {acc_new:.3f}")
\end{lstlisting}

\section{Anexo C: Guía de reproducibilidad}

Esta sección describe de forma precisa el entorno, la secuencia de instalación y los comandos
exactos necesarios para reproducir todos los experimentos presentados en este documento.

\subsection{C.1 Entorno de ejecución}

\textbf{Hardware}: Servidor \texttt{fermidi1-informat-uv-es}, Universitat de València. GPU NVIDIA
disponible en el nodo de cómputo.

\textbf{Software}:

\begin{table}[htbp]
  \centering
  \caption{Versiones del entorno software.}
  \label{tab:software}
  \begin{tabular}{ll}
    \toprule
    \textbf{Componente} & \textbf{Versión} \\
    \midrule
    Python              & 3.10 \\
    PyTorch             & 2.x (con CUDA) \\
    TensorFlow / Keras  & 2.11 (solo para experimentos Keras) \\
    NumPy               & 1.21 \\
    Cython              & 3.0 \\
    MLflow              & 2.x \\
    Plotly              & 5.x \\
    \bottomrule
  \end{tabular}
\end{table}

\subsection{C.2 Instalación del entorno}

\begin{lstlisting}[language=bash, caption={Configuración del entorno desde cero.}]
# 1. Clonar el repositorio del proyecto
git clone <url_repositorio>
cd qsimov

# 2. Crear entorno virtual Python 3.10
python3.10 -m venv venv
source venv/bin/activate

# 3. Instalar dependencias
pip install -r requirements.txt

# 4. Compilar las extensiones Cython (OBLIGATORIO antes del primer uso)
#    Sin este paso, las importaciones del modulo paths/ fallaran.
python setup.py build_ext --inplace
\end{lstlisting}

\subsection{C.3 Descarga y preprocesado de datos}

Cada experimento descarga y preprocesa sus propios datos automáticamente si no se especifica
\texttt{--skip-data-preparation}. Los datasets se almacenan en \texttt{data/} dentro del directorio
del proyecto. Para ImageNet es necesario proporcionar los archivos descargados previamente (ver
\texttt{experiments/imagenet\_subset\_by\_splits/preprocess\_data.py}).

\subsection{C.4 Comandos de ejecución por experimento}

Todos los comandos se ejecutan desde la raíz del directorio \texttt{qsimov/} con el entorno virtual
activo. Los resultados (modelos, historiales, gráficas HTML) se guardan en \texttt{results/} y un
resumen del experimento en \texttt{experiments/<nombre>/run\_exports/}.

\textbf{Experimento 1: Velocidad y precisión (MNIST)}

\begin{lstlisting}[language=bash]
python3 experiments/mnist_speed_loss/main.py \
    --framework pytorch \
    --processor gpu \
    --epochs 10
\end{lstlisting}

Resultados en: \path{results/mnist/learning_rate_speed_loss/pytorch/gpu/}

\textbf{Experimento 2: Robustez al learning rate (MNIST)}

\begin{lstlisting}[language=bash]
python3 experiments/mnist_learning_rate/main.py \
    --framework pytorch \
    --processor gpu \
    --epochs 10
\end{lstlisting}

Resultados en: \path{results/mnist/learning_rate_speed_loss/pytorch/gpu/}

\textbf{Experimento 3: Olvido catastrófico (MNIST)}

\begin{lstlisting}[language=bash]
python3 experiments/mnist_forgetting/main.py \
    --framework pytorch \
    --processor gpu \
    --base-epochs 10 \
    --epochs 5
\end{lstlisting}

Resultados en: \texttt{results/mnist/forgetting/pytorch/gpu/}

\textbf{Experimento 4: Datos escasos y splits (CIFAR-10)}

\begin{lstlisting}[language=bash]
python3 experiments/cifar10_gradient_by_splits/main.py \
    --framework pytorch \
    --processor gpu \
    --model_name lenet \
    --initial-layer -1 \
    --epochs 10 \
    --splits 100 1000 5000
\end{lstlisting}

Resultados en: \texttt{results/cifar10/gradient\_by\_splits/pytorch/gpu/}

\textbf{Experimento 5: Olvido catastrófico (CIFAR-10)}

\begin{lstlisting}[language=bash]
python3 experiments/cifar10_forgetting/main.py \
    --framework pytorch \
    --processor gpu \
    --base-epochs 20 \
    --epochs 5
\end{lstlisting}

Resultados en: \texttt{results/cifar10/forgetting/pytorch/gpu/}

\textbf{Experimento 6: Comparativa de splits (ImageNet-100)}

\begin{lstlisting}[language=bash]
python3 experiments/imagenet_subset_by_splits/main.py \
    --framework pytorch \
    --processor gpu \
    --model-name vgg16 \
    --initial-layer -2 \
    --epochs 10 \
    --splits 1000 5000 10000
\end{lstlisting}

Resultados en: \path{results/imagenet_subset/gradient_by_splits/pytorch/gpu/}

\textbf{Experimento 7: Aprendizaje continual incremental (ImageNet-100)}

\begin{lstlisting}[language=bash]
python3 experiments/imagenet_continual_learning/main.py \
    --framework pytorch \
    --processor gpu \
    --epochs 10
\end{lstlisting}

Resultados en: \path{results/imagenet_subset/continual_learning/pytorch/gpu/}

\textbf{Experimento 8: Barrido de initial\_layer (ImageNet-100)}

\begin{lstlisting}[language=bash]
python3 experiments/initial_layer_sweep/main.py \
    --framework pytorch \
    --processor gpu
\end{lstlisting}

Resultados en: \path{results/imagenet_subset/initial_layer_sweep/pytorch/gpu/}

\textbf{Experimento 9: Streaming incremental de clases (ImageNet-100)}

\begin{lstlisting}[language=bash]
python3 experiments/imagenet_streaming_incremental/main.py \
    --framework pytorch \
    --processor gpu
\end{lstlisting}

Resultados en: \path{results/imagenet_subset/streaming_incremental/pytorch/gpu/}

\subsection{C.5 Consulta de resultados}

Los resultados de cada ejecución se registran automáticamente en MLflow. Para visualizar el panel de
MLflow:

\begin{lstlisting}[language=bash]
mlflow ui --port 5000
\end{lstlisting}

Acceder en el navegador a \texttt{http://localhost:5000}. Cada experimento aparece bajo su nombre
(\texttt{mnist\_speed\_loss}, \texttt{mnist\_learning\_rate}, etc.) con todos los parámetros de
ejecución registrados.

Las gráficas HTML generadas pueden abrirse directamente en cualquier navegador sin servidor
adicional.

\section{Anexo D: Software desarrollado}

Además de la librería Qsimov (preexistente), este TFG ha requerido el desarrollo de un conjunto de
scripts de experimentación. A continuación se describen los módulos principales desarrollados,
agrupados por experimento.

\subsection{D.1 Módulo de infraestructura de experimentos}

\bpath{experiments/mlflow.py}: Utilidades de orquestación
MLflow: \texttt{run\_script()} (ejecuta subprocesos capturando stdout/stderr y registrando el
resultado en MLflow), \texttt{get\_or\_make\_export\_directory()} (crea el directorio
\texttt{run\_exports/<timestamp>\_<run\_name>/}), \texttt{export\_active\_run()} (serializa el
run activo de MLflow a \texttt{report.json}).

\bpath{experiments/git.py}: Comprueba que el repositorio git esté limpio antes de
lanzar un experimento (\texttt{check\_git\_clean()}), y registra el hash del commit activo en
MLflow para garantizar reproducibilidad.

\bpath{experiments/path_utils.py}: Centraliza las rutas de disco de cada experimento
(directorios de resultados, datasets, directorios de exportación). Cualquier cambio en la
organización de carpetas solo requiere modificar este fichero.

\subsection{D.2 Experimento 1: Velocidad y precisión (MNIST)}

\bpath{experiments/mnist_speed_loss/pytorch/train_pytorch_models.py}: Entrena el
modelo base estándar PyTorch sobre MNIST con distintas fracciones del dataset (\texttt{half},
\texttt{quarter}, \texttt{full}) y dos funciones de pérdida (\texttt{categorical\_crossentropy},
\texttt{mse}). Guarda historia de entrenamiento (accuracy, loss, tiempo por epoch) como
\texttt{.pkl}.

\bpath{experiments/mnist_speed_loss/pytorch/train_qsimov_models.py}: Entrena
\texttt{PytorchQsimovGradient} y \texttt{PytorchQsimovLinearSystem} sobre el mismo dataset.
Comparte el PathSelector del modelo base (mismo \texttt{initial\_layer=-2}).

\bpath{experiments/mnist_speed_loss/speed_loss_plots.py}: Genera gráficas HTML
comparando accuracy y loss en función del tiempo de entrenamiento (eje X = tiempo en segundos) para
los tres métodos. Usa Plotly para generar gráficas interactivas.

\bpath{experiments/mnist_speed_loss/main.py}: Orquestador MLflow: lanza los tres
scripts anteriores en orden, registra parámetros y copia los artefactos generados a
\texttt{run\_exports/}.

\subsection{D.3 Experimento 2: Robustez al learning rate (MNIST)}

\bpath{experiments/mnist_learning_rate/train_pytorch_qsimov_models.py}: Itera
sobre 6 valores de learning rate, entrenando para cada uno tanto el modelo estándar PyTorch como
\texttt{PytorchQsimovGradient}. Registra la accuracy final en test para cada combinación.

\bpath{experiments/mnist_learning_rate/learning_rate_plots.py}: Genera gráficas de
accuracy en función del learning rate (escala logarítmica) y curvas de entrenamiento por epoch para
cada LR.

\bpath{experiments/mnist_learning_rate/main.py}: Orquestador MLflow del experimento.

\subsection{D.4 Experimentos de olvido catastrófico (MNIST y CIFAR-10)}

\bpath{experiments/mnist_forgetting/pytorch_train_base_model.py} y
\bpath{experiments/cifar10_forgetting/pytorch_train_base_model.py}: Entrenan el
modelo base sobre la fase 1 (clases 0--4) y guardan el modelo en disco.

\bpath{experiments/mnist_forgetting/pytorch_run_forgetting.py} y
\bpath{experiments/cifar10_forgetting/pytorch_run_forgetting.py}: Ejecutan todos
los métodos de actualización (QsimovLinearSystem acumulado, QsimovGradient, fine-tuning estándar,
re-entrenamiento acumulado) sobre la fase 2 (clases 5--9) y evalúan la accuracy en ambos conjuntos
de clases.

{\tolerance=9999\relax
\bpath{experiments/mnist_forgetting/plot_forgetting.py} y
\bpath{experiments/cifar10_forgetting/plot_forgetting.py}:
Generan gráficas de barras comparando accuracy en clases antiguas vs.\ nuevas para cada método.
\par}

\subsection{D.5 Experimentos de ImageNet}

\bpath{experiments/imagenet_subset_by_splits/train_pytorch.py}: Entrena el
clasificador estándar VGG16 (base convolucional congelada, top ancho 25088$\rightarrow$4096$\rightarrow$4096$\rightarrow$100)
con distintos tamaños de split.

\bpath{experiments/imagenet_subset_by_splits/train_pytorch_qsimov.py}: Entrena el
PathSelector VGG16 (top estrecho 25088$\rightarrow$2048$\rightarrow$128$\rightarrow$100) y
\texttt{PytorchQsimovGradient} para cada split.

\bpath{experiments/imagenet_continual_learning/train_pytorch_qsimov.py}: Implementa
el escenario de 4 rondas de 25 clases nuevas: entrena \texttt{QsimovLinearSystem} con acumulación
de ecuaciones entre rondas, \texttt{QsimovLinearSystem} con reset, y \texttt{QsimovGradient}.

\bpath{experiments/imagenet_streaming_incremental/train_pytorch_streaming.py}:
Simula el escenario de streaming: 20 batches de 5 clases nuevas cada uno. Para cada batch entrena
los 4 métodos comparados y evalúa sobre el test set acumulado.

\bpath{experiments/initial_layer_sweep/train_pytorch_sweep.py}: Prueba distintos
valores de \texttt{initial\_layer} sobre VGG16, midiendo número de caminos generados, tiempo de
construcción del PathSelector, tiempo de entrenamiento y precisión final.


\addcontentsline{toc}{chapter}{Referencias}
\bibliographystyle{unsrt}
\bibliography{bib/bibliografia}

\end{document}
